\section{Μεθοδολογία}

Στο κεφάλαιο αυτό παρουσιάζεται αναλυτικά η μεθοδολογία που χρησιμοποιήθηκε κατά τη διεξαγωγή των πειραμάτων. Συμπεριλαμβάνονται επακριβείς πληροφορίες για την προεπεξεργασία των δεδομένων, τη δομή του μοντέλου και για το post-processing που εφαρμόζεται.

\subsection{Προεπεξεργασία Δεδομένων}

Η προεπεξεργασία των δεδομένων αποτελεί βασικό τμήμα κάθε μελέτης στη Βαθιά Μάθηση. Αποτελεί τη διαδικασία κατά την οποία τα δεδομένα, τα οποία αρχικά βρίσκονται σε μορφή που δεν μπορούν να υποστηρίξουν μάθηση (στην περίπτωσή μας, τα αρχικά βίντεο των ομιλητών με τον ήχο και πριν την απόδοση φωνημάτων σε κάθε frame), αποκτούν τη μορφή που τους επιτρέπει να εισαχθούν στο μοντέλο. Αξιοποιώντας τις γνώσεις πάνω στο αντικείμενο, είναι δυνατόν να μειωθεί ο χρόνος εκπαίδευσης μέσω της μείωσης των διαστάσεων των πληροφοριών που εισάγονται στο νευρωνικό δίκτυο και, γενικότερα, να βελτιωθούν τα τελικά αποτελέσματα. Παρακάτω αναλύονται οι εκάστοτε τεχνικές pre-processing που εφαρμόζονται ώστε να λάβουν τα δεδομένα την επιθυμητή μορφή.

\subsubsection{Εξαγωγή της περιοχής των χειλιών}

Έστω ένα RGB βίντεο ομιλητή $H \times W \times 3 \times T$ με $H$ ύψος και $W$ πλάτος σε pixel και $T$ η χρονική διάσταση (ο αριθμός των frames). Για να περιοριστεί το μέγεθος των δεδομένων και να παραλειφθούν περιττές πληροφορίες του παρασκηνίου αλλά και του προσώπου του ομιλητή, αναπτύχθηκε ένα pre-processing script εντοπισμού της περιοχής των χειλιών. Η περιοχή αυτή είναι που αφορά το αντικείμενο, καθώς εκεί βρίσκεται η πληροφορία των visemes (άλλωστε, για αυτό αποκαλείται lip reading). Χρησιμοποιώντας το διαθέσιμο προεκπαιδευμένο MediaPipe Face Landmarker της Google, εντοπίζονται τα διάφορα ``τοπόσημα προσώπου'' (facial landmarks) του ομιλητή για κάθε frame του αρχικού βίντεο. Οι δείκτες των τοποσήμων των χειλιών εξάγονται για κάθε frame ξεχωριστά και παράγουν μία νέα εικόνα της περιοχής αυτής (συμπεριλαμβανομένου ενός επιπρόσθετου padding γύρω από τα χείλη για ασφάλεια) με διαστάσεις $H' \times W' \times 3$ και $H'<H, W'<W$.

Από τα εκάστοτε frames παράγεται το νέο βίντεο της περιοχής των χειλιών, το οποίο μετασχηματίζεται by default στην τελική διάσταση $224 \times 224 \times 3 \times T$. Κατά τη διάρκεια των πειραμάτων, βρέθηκε ότι ακόμα λιγότερα pixels είναι ικανά να αποδόσουν το περιεχόμενο των visemes, με αποτέλεσμα πριν την εκπαίδευση να εφαρμόζεται ένας ακόμα μετασχηματισμός στα 96 pixels. Κατά τον εντοπισμό των τοποσήμων χρησιμοποιείται και ένας μηχανισμός fallback, σε περίπτωση που το Face Landmarker αδυνατεί να βρει τους ζητούμενους δείκτες. Συγκεκριμένα, εάν δεν εντοπιστούν για ένα δεδομένο frame τα τοπόσημα, τότε επαναχρησιμοποιούνται αυτά του προηγούμενου frame, ενώ αν δεν υπάρχουν ούτε αυτά εκτιμάται η τοποθεσία των χειλιών ως η κάτω περιοχή του προσώπου (50\% του πλάτους και 35\% του ύψους, με κατακόρυφο offset 42\%). Τα δεδομένα κωδικοποιούνται σε μορφή βίντεο .mp4 διατηρώντας το framerate του αρχικού βίντεο, ενώ παράγονται και μερικά παραδείγματα frames από το βίντεο bbaf2n.mpg του dataset GRID (βλ. παράγραφο \ref{Datasets}), τα οποία παρουσιάζονται παρακάτω.

\begin{figure}[H]
	\centering
	\includegraphics[width=0.7\textwidth]{bbaf2n_frame0000_comparison.jpg}
	\caption{Εξαγωγή χειλιών για το frame 0}
\end{figure}

\begin{figure}[H]
	\centering
	\includegraphics[width=0.7\textwidth]{bbaf2n_frame0024_comparison.jpg}
	\caption{Εξαγωγή χειλιών για το frame 24}
\end{figure}

\begin{figure}[H]
	\centering
	\includegraphics[width=0.7\textwidth]{bbaf2n_frame0049_comparison.jpg}
	\caption{Εξαγωγή χειλιών για το frame 49}
\end{figure}

\begin{figure}[H]
	\centering
	\includegraphics[width=0.7\textwidth]{bbaf2n_frame0074_comparison.jpg}
	\caption{Εξαγωγή χειλιών για το frame 74}
\end{figure}

\subsubsection{Εξαγωγή φωνημάτων για κάθε frame}
\label{PhonemeExtraction}
Για να έχει νόημα η εφαρμογή του όρου cross entropy, είναι απαραίτητο να δημιουργηθούν labels για κάθε frame των βίντεο. Αυτό επιτυγχάνεται με συνδυασμό της γνώσης του περιορισμένου περιβάλλοντος και μεθόδων αναγνώρισης φωνημάτων από ηχητικά σήματα ομιλίας. Να σημειωθεί ότι η παρακάτω διαδικασία μπορεί κάλλιστα να εφαρμοστεί και σε μη πεπερασμένο περιβάλλον, αλλά οι περαιτέρω περιορισμοί που προκύπτουν από τέτοια περιβάλλοντα επιτρέπουν απλοποιήσεις και διορθώσεις σε περίπτωση λανθασμένης προφοράς από τον ομιλητή.

Θεωρώντας γνωστό το φραστικό περιεχόμενο κάθε βίντεο (π.χ. στο GRID dataset γίνεται γνωστή η φράση που εκφέρεται από τον ομιλητή μέσω του τίτλου του αρχείου), μπορεί να μετατραπεί η φράση σε αλληλουχία φωνημάτων μέσω του eSpeak-NG, δίχως ακόμα να υπάρχει ωστόσο γνώση για τον χρονισμό του κάθε φωνήματος. Χρησιμοποιείται ελαφριά κανονικοποίηση στα φωνήματα που εξάγονται ώστε το ``αλφάβητο'' των τελικών αλληλουχιών να ταυτίζεται με αυτό του μοντέλου wav2vec2, που χρησιμοποιείται για την εξαγωγή φωνημάτων από το ηχητικό περιεχόμενο. Το CTC μοντέλο του wav2vec2 εξάγει πιθανότητες κάθε φωνήματος για κάθε frame και, συγκρίνοντας με την αλληλουχία φωνημάτων που γνωρίζουμε ότι αναμένεται, παράγει μία τελική πρόβλεψη για την αρχή και το τέλος του κάθε φωνήματος. Μέσω αυτών των πληροφοριών, είναι δυνατόν να ολοκληρωθεί το labelling κάθε frame με το αντίστοιχο φώνημα ή το blank token, εφόσον το frame αντιστοιχεί σε χρονική στιγμή εκτός των φωνημάτων. Έτσι, πέραν της εξαγωγής των απαιτούμενων labels, αποφεύγονται τυχόντα προβλήματα από αστοχία κατά την προφορά της κάθε φράσης.

Η επεξεργασία των φωνημάτων δεν ολοκληρώνεται εκεί. Συγκεκριμένα, δημιουργείται ένα viseme grouping (ομαδοποίηση, δηλαδή, των όμοιων οπτικά φωνημάτων) με βάση τα χαρακτηριστικά των φωνημάτων. Για αναλυτικότερες πληροφορίες ως προς την ομαδοποίηση των φωνημάτων, ανατρέξτε στην ενότητα \ref{Phonemes}. Από τα hard targets που υπάρχουν ήδη για κάθε frame των βίντεο, δημιουργούνται soft targets με βάση τις παραπάνω ομάδες. Συγκεκριμένα, το 85\% αποδίδεται στο πραγματικό φώνημα, ενώ το υπόλοιπο 15\% μοιράζεται στα όμοια φωνήματα του πραγματικού με βάρη που ανταποκρίνονται στην ομοιότητα των φωνημάτων. Για ίδιο viseme, χρησιμοποιείται βάρος 0.95, ενώ για διαφορετικό υπολογίζεται similarity score ως εξής: κάθε φώνημα έχει τον τόπο, τον τρόπο άρθρωσης και το αν είναι έμφωνο ή όχι. Με βάση το πόσα κοινά features άρθρωσης έχουν δύο φωνήματα που δεν ανήκουν στο ίδιο viseme, αποδίδεται μεγαλύτερο βάρος σε αυτά με περισσότερες ομοιότητες και μικρότερο σε αυτά με τις λιγότερες. Έτσι, υπολογίζεται το articulatory\_similarity ως εξής: $sim = 1 - diffs / 3$, όπου $diffs$ είναι ο αριθμός διαφορετικών features. Το τελικό similarity matrix χτίζεται με τον ακόλουθο κανόνα:
$mat[i,j] = \begin{cases}
1 & i = j \\
0 & \text{είτε το }i \text{ είτε το }j \text{ είναι blank token και } i \neq j\\
0.95 & \text{ίδια ομάδα viseme} \\
sim(a,b) \cdot 0.4 & \text{σε άλλη περίπτωση} 
\end{cases}
$

Τελικά τα soft targets χτίζονται με το 85\% στο πραγματικό φώνημα και το υπόλοιπο 15\% μοιρασμένο στα όμοια φωνήματα σε αναλογία με το similarity score των φωνημάτων, εκτός από την περίπτωση του blank token, κατά την οποία διατηρούνται τα hard targets. Αυτό συμβαίνει ώστε να ενισχυθεί η ικανότητα του μοντέλου να προβλέπει ακριβώς πότε λέγεται ένα φώνημα και πότε όχι, με αποτέλεσμα οι χρονισμοί των προβλέψεων να είναι ορθοί.

\subsubsection{Data Augmentation}

Για να εμπλουτιστεί ιδιαίτερα κάποιο περιορισμένο σε μέγεθος dataset, κρίνεται αναγκαίο να εφαρμοστούν μέθοδοι data augmentation, δηλαδή επαύξησης δεδομένων. Αυτές οι μέθοδοι, εκμεταλλευόμενες τη γνώση μας στο αντικείμενο, παράγουν νέα δεδομένα από τα υπάρχοντα και εμπλουτίζουν τη μάθηση του μοντέλου. Είναι σημαντικό να τονιστεί ότι οι τεχνικές data augmentation που εφαρμόζονται πρέπει να συμβαδίζουν με τους στόχους μάθησης του μοντέλου. Για παράδειγμα, μία πιθανή μέθοδος data augmentation στην περίπτωση που τα δεδομένα βρίσκονται σε μορφή βίντεο είναι η αναστροφή της αλληλουχίας των frames. Ωστόσο, όταν το μοντέλο πρόκειται να μάθει σε κάποιο παρακείμενο στο οποίο έχει σημασία η σωστή χρονική αλληλουχία (όπως για παράδειγμα στο lip reading, όπου οι κινήσεις των χειλιών για την παραγωγή ενός φωνήματος είναι ευαίσθητες στη σωστή χρονική σειρά), δε θα μπορούσε να εφαρμοστεί η παραπάνω μέθοδος.

Οι τεχνικές data augmentation που τελικά εφαρμόζονται, λόγω της ορθότητάς τους για τον τελικό στόχο και λόγω της ταχείας εφαρμογής τους σε πολλαπλά δεδομένα, είναι ο καθρεπτισμός γύρω από τον κατακόρυφο άξονα (horizontal flip), η χρωματική αλλαγή (color jitter) και η αφαίρεση/επικάλυψη (cutout/occlusion). Σε κάθε περίπτωση οι μετασχηματισμοί αυτοί εφαρμόζονται σε ολόκληρο το βίντεο και όχι σε μεμονωμένα frames, ώστε να αποφευχθούν ανεπιθύμητα φαινόμενα αστάθειας. Ο καθρεπτισμός εφαρμόζεται με πιθανότητα $p=0.5$, ενώ η χρωματική αλλαγή πραγματοποιείται μέσω των παραγόντων της φωτεινότητας (brightness) με συντελεστή από 0.8 έως 1.2, της αντίθεσης (contrast) με παράγοντα 0.8-1.2 και του κορεσμού (saturation) με παράγοντα 0.9-1.1 και με πιθανότητα $p=0.4$. Τέλος, με τη μέθοδο της επικάλυψης, με πιθανότητα $p=0.25$ επιλέγεται τυχαία μία ορθογώνια περιοχή (8-20\% του ύψους και του πλάτους του frame) και αποκρύπτεται για όλα τα frame του δεδομένου βίντεο, προσομοιώνοντας την ύπαρξη ενός αντικειμένου που εμποδίζει τη θέαση συγκεκριμένης περιοχής. Αυτό έχει ως αποτέλεσμα η εκπαίδευση να γίνεται πιο έυρωστη. Να σημειωθεί ότι η διαδικασία του data augmentation εφαρμόζεται αποκλειστικά στο training split και όχι σε αυτά του validation και του test. Παρακάτω παρουσιάζονται μερικά παραδείγματα των παραπάνω μεθόδων επαύξησης δεδομένων.

\begin{figure}[H]
	\centering
	\includegraphics[width=0.7\textwidth]{horizontal_flip.png}
	\caption{Παρουσίαση του καθρεπτισμού (horizontal flip) για μερικά frames}
\end{figure}

\begin{figure}[H]
	\centering
	\includegraphics[width=0.7\textwidth]{color_jitter.png}
	\caption{Παρουσίαση της αλλαγής χρώματος (color jitter) για μερικά frames}
\end{figure}

\begin{figure}[H]
	\centering
	\includegraphics[width=0.7\textwidth]{random_cutout.png}
	\caption{Παρουσίαση της αφαίρεσης/επικάλυψης (cutout/occlusion) για μερικά frames}
\end{figure}

\subsection{Δομή του προτεινόμενου μοντέλου}

\subsection{Βήματα Post-Processing}

Όπως παρουσιάστηκε το μοντέλο μέχρι τώρα, είναι ικανό μόνο να προβλέψει φωνήματα από το βουβό βίντεο. Στη συγκεκριμένη υποενότητα, περιγράφονται οι τεχνικές post-processing που εφαρμόστηκαν ώστε η αλληλουχία φωνημάτων να μετατραπεί σε λέξεις, βήμα το οποίο οδηγεί στην επίτευξη του αρχικού στόχου του lip reading.

\subsubsection{Αλγοριθμική Πρόβλεψη Λέξεων}

Σε εφαρμογές σε περιορισμένο περιβάλλον, όπου το λεξιλόγιο και, προαιρετικά, η γραμματική των φράσεων είναι γνωστά, είναι δυνατόν η πρόβλεψη των λέξεων από την αλληλουχία φωνημάτων να γίνει αλγοριθμικά και χωρίς τη χρήση τεχνητής νοημοσύνης. Για τα παρακάτω, υποθέτουμε (χωρίς βλάβη της γενικότητας και για να γίνει πιο ξεκάθαρη η επεξήγηση της μεθόδου) ότι η ανάλυση γίνεται στο dataset GRID, όπως παρουσιάζεται στην παράγραφο \ref{Datasets}.

Αρχικά, παράγεται το collapsed sequence φωνημάτων από την αλληλουχία των φωνημάτων που προβλέπεται, ενώνοντας γειτονικά ίδια tokens και απαλείφοντας τα blank tokens, δημιουργώντας μία ομάδα από φωνήματα και τους χρονισμούς τους. Έπειτα, χρησιμοποιείται αλγόριθμος δυναμικού προγραμματισμού ώστε να βρεθεί η τομή των φωνημάτων σε λέξεις και η ανάθεση των προβλεπόμενων λέξεων με το μικρότερο συνολικό κόστος. 

Για κάθε υποψήφια λέξη σε κάθε θέση, υπολογίζεται μία τιμή ομοιότητας με την ακολουθία των φωνημάτων που προβλέπεται. Έστω $O$ η πρόβλεψη και $T$ η υποψήφια λέξη. Ορίζεται η βασική απόσταση ως $D(O,T) = \frac{1}{max(|O|, |T|)} \cdot ED(O,T)$, όπου $ED(O,T)$ είναι η ελάχιστη απόσταση edit της αλληλουχίας φωνημάτων $O$ ώστε να φτάσει στην $T$. Κάθε φώνημα που πρέπει να αλλαχθεί δεν επηρεάζει με ίδιο τρόπο την $ED(O,T)$. Συγκεκριμένα, όμοια φωνήματα (είτε ακουστικά είτε οπτικά) τιμωρούνται λιγότερο σε σχέση με τελείως διαφορετικά φωνήματα. Το κόστος κάθε αλλαγής, από φώνημα $a$ σε φώνημα $b$, υπολογίζεται ως εξής:

$c(a,b) = \begin{cases}
0 & a = b \\
0.15 & \text{ίδια κλάση viseme και ίδιος τρόπος/τόπος εκφοράς, αλλά διαφορετική φω\-νο\-ποί\-η\-ση} \\
0.35 & \text{ίδια κλάση viseme αλλά όχι μόνο διαφορετική φωνοποίηση} \\
min(1, \frac{\Delta (a,b)}{4}) & \text{αλλιώς, με διαφορά ιδιοτήτων $\Delta$} \\
1 & \text{αν δεν υπάρχει καμία ίδια ιδιότητα}
\end{cases}
$

Τότε, το τελικό σκορ ομοιότητας για κάθε υποψήφια λέξη είναι $S(O,T) = D(O,T) - B_{anchor}(O,T) -0.15 B_{rare}(O,T)$, όπου $B_{anchor}$ είναι ένα ισχυρό bonus που δίνεται σε φωνήματα anchor και $B_{rare}$ είναι ένα μικρότερο bonus για σπάνια φωνήματα. Τα φωνήματα anchor είναι φωνήματα που εμφανίζονται μόνο σε μία υποψήφια λέξη κάποιας κατηγορίας. Για παράδειγμα, στις προθέσεις του GRID, το φώνημα /ð/ με ήχο σαν το ελληνικό θήτα εμφανίζεται μόνο στη λέξη with και όχι σε οποιαδήποτε από τις in, by και at. Εάν δεν υπήρχε αυτός ο όρος, μία αλληλουχία φωνημάτων /i/ /ð/ θα οδηγούσε σε πρόβλεψη in και όχι σε πρόβλεψη with, που φαίνεται πιο αληθοφανές, δεδομένου ότι προβλέφθηκε το /ð/ που ανήκει μόνο σε μία λέξη. Με παρόμοιο τρόπο υπολογίζεται και το $B_{rare}$, αλλά πολλαπλασιάζεται με παράγοντα ανάλογο της σπανιότητας του εκάστοτε φωνήματος.

Με την παραπάνω μέθοδο, είναι δυνατόν να προβλεφθούν και λέξεις εντός του περιορισμένου περιβάλλοντος, πέραν των φωνημάτων. Ο αλγόριθμος αυτός σίγουρα μπορεί να αποτελέσει αντικείμενο βελτίωσης και καλύτερης προσαρμογής στο δεδομένο περιβάλλον στο οποίο γίνεται η κάθε μελέτη, ωστόσο κρίθηκε ότι η τωρινή εκδοχή του είναι αρκετή ώστε να επιδείξει ικανοποιητικά αποτελέσματα. Προφανώς, τα παραπάνω δεν μπορούν να εφαρμοστούν σε μη πεπερασμένο περιβάλλον, στο οποίο τόσο το λεξιλόγιο όσο και η γραμματική είναι σαφώς πιο ποικιλόμορφα, καθώς η λογική ιδιαίτερα του δυναμικού προγραμματισμού καταρρέει όσο αυξάνεται η πολυπλοκότητα των πιθανών λέξεων και της τοποθέτησής τους. 